% -----------------------------------------------------------------------
%  lifelines-hc Application Note — Bioinformatics journal
%
%  Class:   oup-authoring-template  (OUP standard; pre-installed on
%           Overleaf and available via CTAN / TeX Live)
%  Style:   contemporary, large, unnumbered section heads, numbered refs
%
%  Compile: pdflatex main && bibtex main && pdflatex main && pdflatex main
%  Overleaf: upload all files; template class is already available there.
% -----------------------------------------------------------------------
\documentclass[unnumsec,webpdf,contemporary,large]{oup-authoring-template}

% ---- packages ----
\usepackage{microtype}
\usepackage{booktabs}
\usepackage{listings}
\usepackage{xcolor}

\lstset{%
  language=Python,
  basicstyle=\ttfamily\small,
  keywordstyle=\color{blue!65!black}\bfseries,
  commentstyle=\color{gray!70},
  stringstyle=\color{green!45!black},
  breaklines=true,
  frame=tb,
  framesep=4pt,
  backgroundcolor=\color{gray!7},
  xleftmargin=1.2em,
  xrightmargin=1.2em,
  aboveskip=0.8em,
  belowskip=0.5em,
  showstringspaces=false,
}

\graphicspath{{../figs/}}

% -----------------------------------------------------------------------
\begin{document}

% ---- journal metadata (leave most blank; editor fills on acceptance) ----
\journaltitle{Biometrika}
\DOI{10.1093/biomet/asaf075}
\copyrightyear{2025}
\pubyear{2026}
\access{Advance Access Publication Date: Day Month Year}
\appnotes{Application Note}

\firstpage{1}

% ---- title ----
\title[{\texttt{lifelines-hc}}: Higher Criticism survival testing]%
      {{\texttt{lifelines-hc}}: a Python package for Higher Criticism
       testing in two-sample survival analysis under non-proportional
       hazards}

% ---- authors ----
\author[1,$\ast$]{Alon Kipnis}
\author[2,$\ast$]{Zohar Yakhini}

\authormark{Kipnis}

\address[1]{%
  \orgdiv{School of Computer Science},
  \orgname{Reichman University},
  \orgaddress{%
    \street{Herzliya},
    \postcode{4610101},
    \country{Israel}%
  }%
}

\corresp[$\ast$]{%
  E-mail: \href{mailto:alon.kipnis@runi.ac.il}{alon.kipnis@runi.ac.il}%
}

% Dates — leave blank; editor fills on acceptance (month arg must be numeric)
\received{}{0}{}
\revised{}{0}{}
\accepted{}{0}{}

%\editor{Associate Editor: Name}

% ---- structured abstract ----
% Bioinformatics Application Notes: Motivation | Results | Availability
% Use \\ between subsections; no blank lines inside \abstract{...}.
\abstract{%
\textbf{Motivation:}
The log-rank test is the standard for two-sample survival comparison and
has good power against a wide range of alternatives. However,
\citet{kipnis_biometrika} showed that it lacks power against
\emph{sparse} hazard departures — cases where the hazard difference is
concentrated in a few time intervals whose locations across the follow-up
are unknown a priori. Weighted log-rank alternatives (Gehan-Wilcoxon,
Tarone-Ware, Peto-Prentice, Fleming-Harrington) each pre-specify a
temporal emphasis, so they too are insensitive to sparse departures
outside their focal window.\\
%
\textbf{Results:}
We introduce \texttt{lifelines-hc}, a Python package extending the
\texttt{lifelines} survival library with the HCHG test — Higher Criticism
(HC) applied to per-interval HyperGeometric (HG) p-values from
\citet{kipnis_biometrika} — for right-censored two-sample survival data.
HCHG detects sparse hazard departures at unknown locations without
committing to any temporal weight. Across three biological case studies
— an immuno-oncology trial (CheckMate~057 PFS, $n=582$), a targeted
kinase-inhibitor trial in metastatic prostate cancer (COMET-1 OS,
$n=1028$), and an adjuvant bisphosphonate trial (AZURE DFS, $n=3359$)
— HCHG achieves $p \leq 0.013$ while the log-rank test and all four
weighted alternatives return $p \geq 0.26$.\\
%
\textbf{Availability and Implementation:}
\texttt{lifelines-hc} is freely available under the MIT licence at
\url{https://github.com/TODO/lifelines-hc} and via
\texttt{pip install lifelines-hc}; it requires Python~${\geq}\,3.8$ and
\texttt{lifelines}~${\geq}\,0.29$. Full analysis scripts and per-method
results tables are provided as Supplementary Data.%
}

\keywords{survival analysis; higher criticism; sparse signal detection;
  hypergeometric test; hypothesis testing; Python}

\maketitle


% =======================================================================
\section{Introduction}
% =======================================================================

The Kaplan-Meier estimator \citep{kaplan1958} and the log-rank test
\citep{mantel1966} are the workhorses of two-sample survival analysis in
biomedical research. The log-rank statistic has good power against a wide
range of alternatives; in particular, \citet{schoenfeld1981} showed it
is asymptotically optimal against proportional-hazards (PH) alternatives.
Nevertheless, \citet{kipnis_biometrika} showed that the log-rank test is
not optimal against \emph{sparse} hazard departures, in which the hazard
difference is concentrated in a small, unknown subset of time intervals.

Such sparse departures arise naturally in several biological settings.
Immune checkpoint inhibitors require an immune-priming phase before
generating durable tumour suppression, producing a crossing-curve pattern
in which the excess hazard is confined to the early and late phases of
follow-up, as in the second-line NSCLC trial CheckMate~057
\citep{borghaei2015}. Targeted therapies whose mechanism is pathway-specific
can produce a transient benefit window that fades as resistance develops:
cabozantinib improved bone scan response at week~12 in the COMET-1 mCRPC
trial \citep{smith2016} but the OS advantage was confined to the
bone-response phase. Benefits conditional on a patient subgroup generate
composite hazard differences concentrated in the period when the susceptible
fraction grows: in the AZURE trial of adjuvant zoledronic acid, the
bone-microenvironment benefit is menopause-dependent and accumulates only
after premenopausal patients transition to postmenopause \citep{coleman2011}.

The need for log-rank alternatives in survival analysis is well recognised,
and several weighted log-rank tests have been developed for non-proportional-hazards
settings \citep{gehan1965,peto1972,taroneware1977,flemingharrington1991}.
Each applies a temporal weight function — emphasising early, intermediate,
or late events — and achieves good power when the signal is located in
the intended window. However, choosing an appropriate weight requires
knowing in advance \emph{where} the hazard departure will occur.
When the departure is sparse but its location is unknown a priori, all
fixed-weight tests can miss it simultaneously, as we demonstrate below.

\citet{kipnis_biometrika} proposed the HCHG test — the Higher Criticism
(HC) statistic of \citet{donoho2004} applied to per-interval
HyperGeometric (HG) p-values — specifically for this setting. HCHG
divides the follow-up into equal-width intervals, derives a hypergeometric
p-value for the event-count comparison in each interval under the null of
equal group-specific hazards, and applies the HC statistic to test whether
the most extreme subset of these p-values is collectively implausible
under the global null. Because HCHG does not pre-specify any temporal
weight, it retains power for sparse departures wherever they occur along
the follow-up axis.

Despite this theoretical advantage, an accessible Python implementation
of HCHG integrated with standard survival analysis workflows has been
lacking. We present \texttt{lifelines-hc}, which fills this gap by
extending the widely-used \texttt{lifelines} library
\citep{davidsonpilon2019}.


% =======================================================================
\section{The \texttt{lifelines-hc} package}
% =======================================================================

\subsection*{The HCHG statistic}

HCHG \citep{kipnis_biometrika} combines two components: hypergeometric
interval p-values and the HC statistic of \citet{donoho2004}.

\textbf{Hypergeometric interval p-values.}
Given two groups with right-censored observations $(T_i, E_i, G_i)$,
partition the follow-up range into $K$ equal-width intervals. For
interval $k$, let $n_k$ be the total events, $r_k^A$ and $r_k^B$ the
subjects at risk in groups A and B, and $x_k$ the events in group B.
Under the null of equal group-specific hazards,
%
\begin{equation}\label{eq:hypergeom}
  x_k \;\Big|\; n_k,\, r_k^A,\, r_k^B
  \;\sim\;
  \mathrm{Hypergeometric}\!\left(r_k^A + r_k^B,\; r_k^B,\; n_k\right),
\end{equation}
%
yielding a one-sided p-value $p_k$ for each interval.

\textbf{HC aggregation.}
Ordering the $K$ p-values as $p_{(1)} \leq \cdots \leq p_{(K)}$, the
HC statistic \citep{donoho2004} is applied to this collection:
%
\begin{equation}\label{eq:hc}
  \mathrm{HCHG}_K
  \;=\;
  \max_{1 \leq j \leq \lfloor\gamma K\rfloor}
  \frac{j/K - p_{(j)}}{\sqrt{p_{(j)}\!\left(1 - p_{(j)}\right)/K}},
\end{equation}
%
where $\gamma \in (0,1)$ (default $0.2$) restricts scanning to the most
extreme fraction of interval p-values \citep{donoho2004}. A global
p-value is obtained by permuting group labels
($n_\mathrm{perms} = 500$ by default). For the two-sided alternative,
the minimum of the two one-sided interval p-values replaces $p_k$
in~(\ref{eq:hc}).

\subsection*{Software design and API}

\texttt{lifelines-hc} exposes a function-based API consistent with
\texttt{lifelines.statistics}, accepting NumPy arrays of event times and
indicators:

\begin{lstlisting}
from lifelines_hc import (higher_criticism_test,
                           fisher_combination_test,
                           min_p_test,
                           suspected_deviations)

result = higher_criticism_test(
    T_A, T_B,
    event_observed_A=E_A, event_observed_B=E_B,
    n_intervals_to_pool=100, n_permutations=500,
    alternative='both', gamma=0.2, seed=42)

print(result.p_value)       # permutation-calibrated p-value
print(result.test_statistic)  # HCHG statistic value
\end{lstlisting}

\noindent
The companion \texttt{suspected\_deviations()} function returns a
\texttt{pandas.DataFrame} of per-interval p-values with a
\texttt{suspected} boolean column marking HCHG-flagged intervals,
enabling overlay of diagnostic shading on Kaplan-Meier plots. Two
complementary omnibus tests are also provided: Fisher combination
\citep{fisher1932}, which aggregates evidence via $-2\sum_k \log p_k$,
and a Bonferroni-corrected MinP test. All three share the same
permutation infrastructure and the same hypergeometric interval p-values.
The package additionally ships with plotting utilities
(\texttt{plot\_km\_with\_hc}, \texttt{plot\_pvalue\_profile}) and a
convenience wrapper \texttt{run\_all\_tests()} that benchmarks HCHG,
Fisher, MinP, log-rank, and the four weighted alternatives in a single
call, returning a \texttt{pandas.DataFrame} for easy comparison.


% =======================================================================
\section{Results}
% =======================================================================

We demonstrate \texttt{lifelines-hc} on three biological datasets with
distinct non-PH structures (Figure~\ref{fig:main}), benchmarking HC
against the log-rank test and four weighted alternatives.
Full per-method p-value tables are provided in Supplementary Data.

\subsection*{Immuno-oncology: crossing survival curves}

CheckMate~057 \citep{borghaei2015} randomised 582 patients with advanced
non-squamous NSCLC to nivolumab ($n=292$) or docetaxel ($n=290$).
Progression-free survival showed crossing curves: docetaxel provided a
short-term advantage during immune priming, subsequently overtaken by the
durable nivolumab response (Figure~\ref{fig:main}A). Individual patient
data were reconstructed from the published Kaplan-Meier curve via the
algorithm of \citet{guyot2012} as implemented in the \texttt{kmdata}
R package.

The log-rank test is non-significant ($p = 0.351$) as the early excess
and late benefit partially cancel. Among the weighted alternatives,
Gehan-Wilcoxon ($p = 0.041$) and Peto-Prentice ($p = 0.049$) reach
borderline significance by detecting the early crossing, but
Tarone-Ware ($p = 0.358$) and Fleming-Harrington\,(1,1) ($p = 0.256$)
do not. HCHG achieves $p = 0.002$ and characterises the full crossing
structure with far greater power (Figure~\ref{fig:main}D) by flagging
both the early divergence and late separation windows.

\subsection*{Targeted therapy in prostate cancer: a transient early benefit}

The COMET-1 trial \citep{smith2016} randomised 1028 patients with
chemotherapy-pretreated mCRPC to cabozantinib ($n=682$) or prednisone
($n=346$) in a 2:1 ratio. Overall survival showed no significant
log-rank difference ($p = 0.262$). Cabozantinib, by inhibiting
MET and VEGFR2 kinases, significantly improved bone scan response at
week~12 (a pre-specified secondary endpoint), confirming genuine
biological activity against bone-metastatic disease.
This short-term activity did not translate into sustained OS benefit,
however, as resistance through alternative survival pathways rapidly
emerged. The result is a temporally concentrated early OS advantage
(months~3--7) that is subsequently lost (Figure~\ref{fig:main}B,E).

All four weighted alternatives are non-significant ($p \geq 0.19$):
the early-weighted Gehan-Wilcoxon ($p = 0.193$) and Peto-Prentice
($p = 0.206$) come closest but are diluted by the null period after
the bone-response window. HCHG ($p = 0.012$) and Fisher combination
($p = 0.020$) aggregate evidence from the specific early intervals
where the transient effect is concentrated.

\subsection*{Adjuvant bisphosphonate therapy: a menopause-dependent effect}

The AZURE trial \citep{coleman2011,coleman2014} randomised 3359 early-stage
breast cancer patients to standard adjuvant therapy with ($n=1681$) or
without ($n=1678$) zoledronic acid. Disease-free survival showed no
significant log-rank difference ($p = 0.305$) because the drug's
bone-microenvironment benefit depends on the low-oestrogen state of
postmenopause, which was absent in the many premenopausal patients
enrolled at baseline. As these patients progressively transitioned to
postmenopause during the decade-long follow-up, the hazard difference
accumulated in a specific mid-to-late window (months 20--60) rather than
uniformly (Figure~\ref{fig:main}C,F).

All four weighted alternatives are non-significant ($p \geq 0.28$):
no single temporal weight captures a pattern distributed across this
window while remaining unaffected by the null early period. HC
($p = 0.012$) detects this scattered but collectively significant signal.
The result is validated by published subgroup analyses showing significant
DFS benefit in postmenopausal patients \citep{coleman2011,coleman2014},
confirming that HCHG is detecting a genuine biological effect.

Across the three case studies, the two borderline significances for
Gehan-Wilcoxon and Peto-Prentice in CheckMate~057 are the only instances
where a weighted alternative reaches $p < 0.05$; these reflect their
early emphasis coincidentally aligning with the crossing-curve signal
rather than general sensitivity to non-PH. Fisher combination achieves
$p \leq 0.04$ in all three domains, consistent with the sparse, localised
nature of the signals.


% =======================================================================
\section{Conclusion}
% =======================================================================

\texttt{lifelines-hc} provides an accessible Python implementation of
HCHG for two-sample right-censored survival data. HCHG is specifically
designed for sparse hazard departures at unknown temporal locations:
unlike weighted log-rank alternatives, it requires no a priori
specification of where the hazard difference will occur and retains
power for any concentrated departure along the follow-up axis. The
package integrates transparently with \texttt{lifelines} workflows,
returns diagnostic interval-level information identifying which time
windows drive the signal, and includes Fisher combination and MinP tests
as complementary omnibus alternatives sharing the same hypergeometric
infrastructure.


% =======================================================================
%  Back matter
% =======================================================================

\section*{Acknowledgements}
TODO.

\section*{Conflict of Interest}
None declared.

\section*{Funding}
The work of Alon Kipnis was supported in parts by the US-Israel Binational Science Foundation (BSF) grant 2022124.

\section*{Data availability}
Individual patient data for all three trials (CheckMate~057, COMET-1,
and AZURE) were reconstructed from published Kaplan-Meier curves using
the algorithm of \citet{guyot2012} as implemented in the \texttt{kmdata}
R package (\url{https://github.com/raredd/kmdata}).
All analysis code is available at
\url{https://github.com/TODO/lifelines-hc/tree/main/AppNote}.


% =======================================================================
%  Figure — full-width across both columns
% =======================================================================

\begin{figure*}[t!]
  \centering
  \includegraphics[width=\textwidth]{figure1}
  \caption{%
    \textbf{Higher Criticism detects non-proportional hazard departures
    across three biological domains with distinct temporal patterns.}
    Top row: Kaplan-Meier curves with HC-flagged intervals (shaded).
    Bottom row: per-interval $-\!\log_{10}$(hypergeometric p-value);
    bars exceeding the HC threshold (dashed, permutation-calibrated)
    are highlighted in red.
    %
    \textbf{(A,D)}~CheckMate~057 PFS (nivolumab vs.\ docetaxel, $n=582$);
    log-rank $p=0.351$, HCHG $p=0.002$. Curves cross at $\sim\!3$ months
    (immune-priming excess) before nivolumab gains a durable advantage.
    %
    \textbf{(B,E)}~COMET-1 OS (cabozantinib vs.\ prednisone, $n=1028$);
    log-rank $p=0.262$, HCHG $p=0.012$. Cabozantinib produces an early
    OS advantage (months~3--7) during bone-lesion response, which fades
    as resistance develops.
    %
    \textbf{(C,F)}~AZURE DFS (zoledronic acid vs.\ control, $n=3359$);
    log-rank $p=0.305$, HCHG $p=0.012$. The signal reflects a
    menopause-dependent benefit accumulating in months~20--60.
    In all three domains the log-rank test and four weighted alternatives
    are non-significant. Individual patient data reconstructed
    via \citet{guyot2012}.
  }
  \label{fig:main}
\end{figure*}


% =======================================================================
%  Bibliography
% =======================================================================

% oup-authoring-template.bst is available on Overleaf; locally use plainnat.
% Switch to \bibliographystyle{oup-authoring-template} when uploading to Overleaf.
\bibliographystyle{plainnat}
\bibliography{refs}

\end{document}
